\section{Понятие аналитической функции в точке и в области. Понятие
гармонической функции. Связь между аналитическими и гармоническими
функциями.}\label{ux43fux43eux43dux44fux442ux438ux435-ux430ux43dux430ux43bux438ux442ux438ux447ux435ux441ux43aux43eux439-ux444ux443ux43dux43aux446ux438ux438-ux432-ux442ux43eux447ux43aux435-ux438-ux432-ux43eux431ux43bux430ux441ux442ux438.-ux43fux43eux43dux44fux442ux438ux435-ux433ux430ux440ux43cux43eux43dux438ux447ux435ux441ux43aux43eux439-ux444ux443ux43dux43aux446ux438ux438.-ux441ux432ux44fux437ux44c-ux43cux435ux436ux434ux443-ux430ux43dux430ux43bux438ux442ux438ux447ux435ux441ux43aux438ux43cux438-ux438-ux433ux430ux440ux43cux43eux43dux438ux447ux435ux441ux43aux438ux43cux438-ux444ux443ux43dux43aux446ux438ux44fux43cux438.}

\textbf{Определение:} Функция \(f(z)\) \emph{аналитическая}:

\begin{enumerate}
\def\labelenumi{\arabic{enumi}.}
\tightlist
\item
  в области \(G\), если \(f'(z)\) непрерывна в \(G\) (т.е.
  \(\forall z\in G \ \exists f'\))
\item
  в точке \(z=z_{0}\), если она аналитическая в \(U_{\delta}(z_{0})\)
\item
  в точке \(z=\infty\), если она аналитическая в \(U(\infty)\).
\end{enumerate}

\textbf{Определение:} Функция \(u(x,y)\) \emph{гармоническая} в
\(G\subset \mathbb{R}^{2}\), если \(u(x,y)\in C^{2}(G)\) и удовлетворяет
уравнению Лапласа: \[
\frac{\partial^{2}u}{\partial x^{2}}+ \frac{\partial^{2}u}{\partial y^{2}}\equiv 0\text{ в }G
\]

\textbf{Теорема 1:} Пусть \(f(z)\) аналитическая в \(G\),
\(f(z)=u(x,y)+iv(x,y)\), \(u(x,y),v(x,y)\in C^{2}(G)\). Тогда \(u(x,y)\)
и \(v(x,y)\) - гармонические функции.

\textbf{Теорема 2:} Пусть дана функция \(u(x,y)\) гармоническая в
\(\mathcal{D}\) и \(\mathcal{D}\) - односвязная область. Тогда
\(\exists v(x,y)\in C^{2}(\mathcal{D})\) такая, что функция
\(f(z)=u(x,y)+iv(x,y), \ z=x+iy\) - аналитическая функция в
\(\mathcal{D}\).
